\section{Processi Primari}
\subsection{Fornitura}
Il processo di fornitura regola i rapporti del gruppo con il proponente e con il committente. All'interno del progetto didattico:
\begin{itemize}
	\item Il \textbf{docente} ricopre il ruolo di committente rispetto alle forniture, alle scadenze, ai vincoli regolamentari e alla valutazione;
	\item Il gruppo \nomegruppo ricopre il ruolo di fornitore, assumendosi la responsabilità di produrre, verificare e consegnare la documentazione richiesta dalle revisioni di avanzamento.
\end{itemize}

Il processo comprende:
\begin{itemize}
	\item La valutazione dei capitolati;
	\item La candidatura al capitolato selezionato;
	\item La formalizzazione degli impegni economici e temporali;
	\item La produzione della documentazione richiesta;
	\item La gestione delle comunicazioni ufficiali
\end{itemize}

\subsubsection{Strumenti a supporto}
Il gruppo si avvale dei seguenti strumenti:
\begin{itemize}
	\item \textbf{Email istituzionale del gruppo}: l'indirizzo \emailgruppo viene utilizzata per tutte le comunicazioni ufficiali con committente e proponente.
	\item \textbf{Repository GitHub}: essa sarà il punto di raccolta, versionamento e pubblicazione di tutta la documentazione e del codice del POC;
	\item \textbf{GitHub Pages}: utilizzato per esporre la documentazione prodotta;
	\item \textbf{GitHub Projects}: per il tracciamento dei TODO, delle Milestone, RTB e PB;
\end{itemize}

\subsubsection{Documentazione prodotta in fase di Candidatura}
Durante la fase di candidatura, il gruppo si è occupato di produrre e approvare tutta la documentazione necessarie per la presentazione della propria proposta.

\begin{table}[htbp]
	\centering
	\caption{Registro documentazione prodotta}
	\label{tab:documenti}
	\renewcommand{\arraystretch}{1.3}
	\small
	\begin{tabularx}{\textwidth}{|>{\centering\arraybackslash}p{4.5cm}
			|X|}
		\hline
		\rowcolor[gray]{0.88}
		\textbf{Documento} &
		\centering\arraybackslash\textbf{Descrizione} \\
		\hline
		% ----- Esempio di compilazione (cancella o modifica) -----
		% Manuale utente & Guida all'utilizzo del sistema per l'utente finale \\ \hline
		% ---------------------------------------------------------
		Lettera di Presentazione & Documento con cui il gruppo comunica in maniera formale la candidatura al capitolato selezionato. Contiene le informazioni sul capitolato, della proponente, della scadenza fissata, del budget preventivato e gli impegni assunti. \\ \hline
		Valutazione dei Capitolati & Documento che contiene il confronto tra tutti i capitolati presentati, analizzati da parte del gruppo secondo aspetti positivi, criticità e rischi. \\ \hline
		Preventivo Costi e Assunzione Impegni & Documento che formalizza il preventivo economico, fissa le scadenze, la distribuzione delle ore produttive e la ripartizione dei ruoli.  \\ \hline
		Norme di Progetto & Documento interno contenete le regole operative. \\ \hline
		Verbali interni & Documento interno il cui scopo è tracciare le decisioni organizzative e le assegnazioni delle attività del gruppo. \\ \hline
		Verbali esterni & Documento esterno che formalizza gli incontri con le aziende, gli argomenti discussi e i chiarimenti ottenuti. \\ \hline
	\end{tabularx}
\end{table}